\section{Overview}
\label{sec:overview}

Every day, many organizations deal with large numbers of documents like invoices, registration forms, and official records. Someone needs to read each one and pull out specific pieces of information, such as a total amount, a name, or a date. Doing this by hand is very slow and expensive. So many organizations now use AI to do it automatically. This is called \textit{data extraction}.

Right now, most organizations do this by sending their documents to an online AI service. The service reads the documents and sends back the information, but it charges a small fee every time. When you have thousands of documents every day, those fees add up very fast. There is also a second problem: private documents have to leave the organization and go to someone else's server, which raises serious privacy and security concerns.

A newer option is now available. Small AI models that run entirely on a laptop with no internet connection---such as Google's Gemma~4 E4B \cite{gemma4_2026}, Microsoft's Phi-4 Mini \cite{phi4mini_2025}, and Alibaba's Qwen2.5-VL 3B \cite{qwen25vl_2025}---are free to use, keep all data private, and can read and understand document images. However, organizations do not yet have clear, honest numbers that show whether these local models are accurate enough, fast enough, or cheap enough compared to the paid online option. There is also an open question about whether the results change depending on how the instruction given to the AI is written.

This proposal describes a study that will test and compare these two approaches side by side. On one side are three local AI models; on the other is a paid online AI service. The comparison covers three things: how accurate each one is (tested using different ways of writing the instruction, to make the comparison fair), how fast each one is, and how much each one costs. The goal is to give organizations clear numbers so they can decide which option is right for them.
